\section{What a group is}\label{sec:def}

\subsection{Multiplying symmetries}

Symmetries of the cube can be combined: do one, then do another, and the result is again a
symmetry. That operation is the multiplication we are going to axiomatise, and it needs a
convention, because ``do $g$ then $h$'' and ``do $h$ then $g$'' are different motions.

\begin{defi}\label{def:order}
For symmetries $g$ and $h$, the product $gh$ means: \emph{first do $h$, then do $g$}.
\end{defi}

The order looks backwards until you notice where it comes from. A symmetry is a function on the
points of space, and functions are written to the left of their inputs: $(g\circ h)(x)=g(h(x))$, so
$h$ acts first. Every computation in this chapter depends on this convention, so it is worth
saying out loud once: \emph{the right-hand factor acts first}.

\sfig{The convention in a picture: the product of two symmetries is the single symmetry that does
the same job as the two in succession.}{\figcompose}

\subsection{The definition}

\begin{defi}\label{def:group}
A \defn{multiplication} on a set $G$ is a function that takes an ordered pair of elements of $G$ and
returns an element of $G$, written $(g,h)\mapsto gh$. A \defn{group} is a set $G$ together with a
multiplication on it such that:
\begin{itemize}[nosep,leftmargin=1.4em]
\item[\textbf{(G1)}] \emph{Identity.} There is an element $1\in G$ with $1g=g1=g$ for every
$g\in G$.
\item[\textbf{(G2)}] \emph{Inverses.} For every $g\in G$ there is an $h\in G$ with $gh=hg=1$. It is
written $g\inv$.
\item[\textbf{(G3)}] \emph{Associativity.} $(fg)h=f(gh)$ for all $f,g,h\in G$.
\end{itemize}
\end{defi}

Two things about the definition are easy to read past.

The first is that the \emph{multiplication is part of the data}. The same set can be a group with
one rule and fail to be one with another: the integers with addition form a group, the integers
with subtraction do not. So a group is a pair, a set and a rule, and naming only the set is an
abbreviation.

The second is that the word ``multiplication'' is a generic label. For the integers it will mean
addition; for symmetries it means composition; for words it will mean writing one after the other.
Nothing about it needs to resemble multiplying numbers. In particular \textbf{nothing in
Definition~\ref{def:group} says $gh=hg$}. Groups in which that happens to hold for every pair are
called \defn{abelian}, and most of the interesting groups in this chapter are not.

\subsubsection*{The rule is a function, and ``multiplication'' is only its name}

Definition~\ref{def:group} says the multiplication is a \emph{function}, and it means that
literally: a function from $G\times G$ to $G$. A function of that shape is called a \defn{binary
operation} on $G$, and the phrase carries two facts that the bare word ``function'' does not. The
inputs are an ordered pair of elements of $G$, and the output lands back in $G$. That second fact
is closure, and it is the one most often violated by a natural-looking candidate.

So the rule can be anything at all, provided it obeys the three axioms. Here is what it actually is
in the examples of this chapter.

\begin{center}
\small
\begin{tabular}{@{}lll@{}}
\toprule
group & the operation is & written\\
\midrule
cube rotations, $D_n$, $S_n$ & composition of functions & $gh$\\
$\Z$, $\Z^2$ & addition & $g+h$\\
$\Z/n\Z$ & add, then take the remainder & $g+h$\\
$\GL(n,\R)$ & matrix multiplication & $gh$\\
$F_2$ & write one string after the other, then cancel & $gh$\\
\bottomrule
\end{tabular}
\end{center}

Only the matrix row is multiplication in the ordinary sense. Concatenating two strings and
cancelling is about as far from multiplying numbers as an operation can get, and it is still a
group.

The name survives for three reasons. Writing the operation as juxtaposition, $gh$, is as compact as
notation gets, and it hands you $1$ for the identity, $g\inv$ for the inverse and $g^n$ for
repetition without inventing anything. The first groups ever studied were substitutions and
numbers, where the operation genuinely was composition or multiplication, and the notation outlived
those examples. And a group carries only \emph{one} operation, so there is no addition alongside it
to be confused with. A set with two compatible operations is a richer object, a ring or a field,
and that is a different subject.

One convention is worth knowing before you meet it. Additive notation, with $0$ for the identity
and $-g$ for the inverse, is reserved almost universally for abelian groups; nobody writes a
non-commutative operation with a plus sign. So the integers get $+$ and the symmetric group gets
juxtaposition, and that choice quietly tells the reader the group is commutative before any axiom
has been checked.

\sfig{The anatomy of the definition: a set, a rule on that set whose answers stay inside the set,
and three promises about the rule.}{\figXj}

\begin{rem}
The word \emph{function} turns up at two levels here, and conflating them causes trouble. The
\emph{operation} is a function of two variables, taking $(g,h)$ to $gh$. In a symmetry group the
\emph{elements} are themselves functions, each one a map from the object to itself. When both
happen at once, the operation is composition of the elements, and that is exactly why associativity
costs nothing for the cube, for $D_n$ and for $S_n$: composing functions is associative whatever
the functions are.
\end{rem}

\begin{rem}
The requirement that $gh$ lie back in $G$ is often called \defn{closure}. In
Definition~\ref{def:group} it is not a separate axiom, because it is built into the word
``function into $G$''. When you are handed a candidate group, however, closure is usually the
first thing to check and the easiest to get wrong; Section~\ref{sec:checklist} has an example
where it fails inside the cube.
\end{rem}

\begin{rem}
The symbol $1$ in (G1) is a \emph{name}, not the number one. For the integers under addition the
identity is $0$; for the cube it is ``leave it alone''; for invertible matrices it is the identity
matrix. Many books write $e$ instead, from the German \emph{Einselement}. Likewise $g\inv$ is a
name for the element that undoes $g$, and it is not ``$1/g$'' unless the group happens to consist
of numbers under ordinary multiplication.
\end{rem}

\begin{rem}
Insisting that the rule be a function is not a formality. It demands that every ordered pair gets
an answer, and that each pair gets exactly \emph{one} answer, so that the result cannot depend on
how you carry the rule out. That second demand is where real work goes twice in this chapter. In
the free group the rule is ``write one word after the other, then cancel'', and the cancelling can
be done in several orders; until every order is shown to reach the same reduced word, the recipe is
not a function and there is no group (Section~\ref{sec:free}). The same demand bites again for
quotients, where the rule on cosets looks fine until you notice the answer can depend on which
representative you pick, and becomes a function only when the subgroup is normal
(Section~\ref{sec:quot}).
\end{rem}

\begin{tryit}
To see how arbitrary the rule may be, define $x*y=x+y+1$ on the real numbers. Check that it is
associative, find its identity and the inverse of $x$, and conclude that it is a group. It is
neither addition nor multiplication, and yet it is a relabelled copy of $\R$ under addition.
\end{tryit}

\subsection{The three axioms on the cube}

The cube is the example to hold in mind, so here is each axiom checked on it. Grey letters are
fixed positions, and a coloured disc names the corner of the original cube now sitting there.

\subsubsection*{(G1) Identity}

``Leave the cube alone'' is a rotation, by $0^\circ$. Doing it after any move $r$ leaves the
picture of $r$ untouched, and doing it before $r$ produces $r$ as well, so $1r=r1=r$.

\sfigs{0.80}{Identity on the cube: start, after a quarter turn $r$, then after doing nothing. The last two
pictures are identical.}{\figXl}

\subsubsection*{(G2) Inverses}

The inverse of a quarter turn is the quarter turn the other way about the same axis. Three more
turns in the original direction do the same job, so $r\inv=r^3$. A half turn is its own inverse,
which is the analogue of $0$ being its own inverse in the integers.

\sfigs{0.80}{Inverses on the cube: start, after $r$, then after $r\inv$. Every disc is back on its grey
letter, so $r\inv r=rr\inv=1$.}{\figXm}

\subsubsection*{(G3) Associativity}

Let $h$ be a quarter turn about the vertical axis, $g$ a quarter turn about the left--right axis,
and $f$ a half turn about the vertical axis. Performed one after another they do this:

\sfigs{0.80}{The three moves $h$, then $g$, then $f$, one picture at a time.}{\figXn}

Now $(fg)h$ says: do $h$, then the \emph{single} rotation that $f$ and $g$ amount to together.
And $f(gh)$ says: do the single rotation that $g$ and $h$ amount to, then $f$. Both routes end at
the same cube as the row above.

\sfigs{0.80}{$(fg)h$: after $h$, the single move $fg$ turns out to be the half turn about the edge axis
through the midpoints of $AB$ and $HG$.}{\figXo}

\sfigs{0.80}{$f(gh)$: the single move $gh$ is a third of a turn about the diagonal $CE$, and then
$f$. Same final cube.}{\figXp}

The brackets only decide which two moves you merge into one before performing anything; the cube
still receives $h$, then $g$, then $f$, so the final position cannot depend on them. Said in one
line: symmetries are functions, and composition of functions satisfies
$\big((fg)h\big)(x)=f\big(g(h(x))\big)=\big(f(gh)\big)(x)$ for every point $x$.

\begin{bigpic}
This is worth remembering as a labour-saving device. \emph{Any} group whose elements are functions
and whose multiplication is composition satisfies (G3) automatically. Associativity only costs
something when the multiplication is given by a formula with cases, as for the integers mod $n$ in
Section~\ref{sec:modn}, or by a recipe on strings, as for the free group in
Section~\ref{sec:free}.
\end{bigpic}

\subsection{Why exactly these three axioms}

The three promises are not a random selection. They are exactly what is needed to solve equations.
Suppose $g$ and $k$ are known and you want $x$ with $gx=k$; on the cube, ``which move $x$, followed
by $g$, produces $k$?''

\begin{center}
\small
\begin{tabular}{@{}ll@{}}
$gx=k$ & given\\
$g\inv(gx)=g\inv k$ & multiply both sides on the left by $g\inv$ \hfill (G2)\\
$(g\inv g)x=g\inv k$ & regroup \hfill (G3)\\
$1x=g\inv k$ & $g\inv g=1$ \hfill (G2)\\
$x=g\inv k$ & \hfill (G1)\\
\end{tabular}
\end{center}

Each axiom is used exactly once, and the solution is unique, since the calculation derives it.
Drop any one axiom and the argument breaks. The same three steps prove the
\defn{cancellation law}: if $gx=gy$ then $x=y$, in every group.

\begin{tryit}
Solve $gx=k$ on the cube with $g$ the quarter turn about the vertical axis and $k$ the half turn
about the same axis. Then solve $xg=k$ for the same $g$ and $k$, and notice that the answer requires
multiplying on the \emph{right}, giving $x=kg\inv$. In an abelian group these two answers agree; in
general they do not.
\end{tryit}

\begin{prop}
In any group the identity is unique, inverses are unique, and $(gh)\inv=h\inv g\inv$.
\end{prop}

\begin{proof}
If $1$ and $1'$ both satisfy (G1) then $1=1\cdot 1'=1'$. If $h$ and $h'$ are both inverses of $g$
then $h=h1=h(gh')=(hg)h'=1h'=h'$. For the last claim, multiply $gh$ by $h\inv g\inv$ and regroup:
$(gh)(h\inv g\inv)=g(hh\inv)g\inv=gg\inv=1$, and similarly on the other side.
\end{proof}

The formula $(gh)\inv=h\inv g\inv$ is the \defn{socks-and-shoes rule}: to undo ``socks, then
shoes'' you take the shoes off first. It will be used constantly from Section~\ref{sec:free}
onwards.

\subsection{Which of these are groups?}\label{sec:checklist}

The fastest way to get used to Definition~\ref{def:group} is to watch it fail.

\begin{center}
\small
\begin{tabular}{@{}llcccc@{}}
\toprule
set & rule & closed & identity & inverses & associative\\
\midrule
$\Z$ & $+$ & \cmark & $0$ & $-n$ & \cmark\\
$\{0,1,2,\dots\}$ & $+$ & \cmark & $0$ & \xmark\ no $-5$ & \cmark\\
$\Z$ & $-$ & \cmark & \xmark\ $0-5\neq5$ & --- & \xmark\\
$\Z$ & $\times$ & \cmark & $1$ & \xmark\ no $\tfrac12$ & \cmark\\
$\{1,-1\}$ & $\times$ & \cmark & $1$ & each its own & \cmark\\
$\R$ & $\times$ & \cmark & $1$ & \xmark\ none for $0$ & \cmark\\
$\R\setminus\{0\}$ & $\times$ & \cmark & $1$ & $1/x$ & \cmark\\
the $24$ rotations of a cube & do one, then the other & \cmark & do nothing & undo & \cmark\\
the $6$ quarter turns of a cube & do one, then the other & \xmark & --- & --- & ---\\
$\{1\}\cup\{\text{3 face half turns}\}$ & do one, then the other & \cmark & $1$ & each its own & \cmark\\
\bottomrule
\end{tabular}
\end{center}

Four of these ten are groups. The failures repay study.
For $\Z$ with subtraction, $5-0=5$ but $0-5\neq 5$, so $0$ is an identity on one side only, and
$(8-4)-2=2$ while $8-(4-2)=6$, so (G3) fails too.
For the non-negative integers, the only problem is inverses; everything else is fine, which shows
the axioms are genuinely independent of one another.
The last line but one is the interesting failure, and it happens inside the cube.

\sfig{Closure can fail inside the cube: a quarter turn followed by a quarter turn about a different
axis is a third of a turn about a long diagonal, not a quarter turn. So the six quarter turns are
not closed under the multiplication, and they do not form a group on their own.}{\figXk}

\begin{why}
Associativity never fails for a set of symmetries, so it is not worth checking there; it fails only
for rules invented by formula, like subtraction. This is why the checklist above has no
\xmark\ in its last column except on the row for subtraction.
\end{why}

\begin{tryit}
Decide, using the columns of the table: the even integers with $+$; the odd integers with $+$; the
eight third-of-a-turn rotations of the cube together with the identity (try two of them about
different diagonals); the four rotations about a single face axis.
\end{tryit}
